
\documentclass[12pt,a4paper]{report}

\setcounter{secnumdepth}{3} %The level of sections that get numbered
\setcounter{tocdepth}{3} %The level of sections to appear in ToC
\usepackage{url}
\usepackage{graphicx}
\usepackage{float}
\usepackage{subcaption}
\usepackage{braket}
\usepackage{calligra}
\usepackage{listings}
\usepackage{hyperref}
\usepackage{graphicx} % Required for inserting images
\usepackage{amssymb}
\usepackage{amsmath}
\usepackage{xcolor}
\usepackage{aurical}
\usepackage[T1]{fontenc}

\definecolor{codegreen}{rgb}{0.133,0.545,0.133}
\definecolor{codegray}{rgb}{0.5,0.5,0.5}
\definecolor{codepurple}{rgb}{0.58,0,0.827}
\definecolor{codeblue}{rgb}{0.0,0.39,0.75}
\definecolor{codeback}{rgb}{0.97,0.97,0.97}

\lstdefinestyle{python}{
  language=Python,
  backgroundcolor=\color{codeback},
  commentstyle=\color{codegreen},
  keywordstyle=\color{codeblue}\bfseries,
  stringstyle=\color{codepurple},
  numberstyle=\tiny\color{codegray},
  basicstyle=\ttfamily\small,
  breakatwhitespace=false,
  breaklines=true,
  keepspaces=true,
  numbers=left,
  numbersep=5pt,
  showspaces=false,
  showstringspaces=false,
  showtabs=false,
  tabsize=4,
  xleftmargin=0pt,
  xrightmargin=0pt,
  framesep=4pt,
  frame=single,
  rulecolor=\color{codegray},
}

\begin{document}
\section{The Forward Model}
The problem is framed as an inverse problem. The forward map takes position and orientation data and calculated 8 flux values from each of the coils. Our goal is to approximate an inverse to this map. To do so, we first try to understand the forward map 
\subsection{Definition and properties.}
First we define a pose vector $\mathbf{p} \in \mathbb{R}^5$ \\
\begin{equation}
  \mathbf{p} = 
  \begin{bmatrix}
    x \\ y \\ z \\ \theta \\ \phi 
  \end{bmatrix}
\end{equation}

In actuality, $\mathbf{p}$ is restricted to a set $\Omega \subset \mathbb{R}^5$

\begin{equation}
  \Omega = \left\{
  \begin{bmatrix} x \\ y \\ z \\ \theta \\ \phi \end{bmatrix} \in \mathbb{R}^5 : 
  \begin{matrix} x \in (-0.125,0.125) \\ y \in (-0.125,0.125) \\ z \in (0,0.25) \\ \theta \in (0,\pi) \\ \phi \in [0,2\pi) \end{matrix}
  \right\}
\end{equation}

Then we define a vector of fluxes, $\mathbf{\Phi} \in \mathbb{R}^8$
\begin{equation}
  \mathbf{\Phi} = 
  \begin{bmatrix}
    \Phi_1 \\ \Phi_2 \\ \Phi_3 \\ \Phi_4 \\ \Phi_5 \\ \Phi_6 \\ \Phi_7 \\ \Phi_8 
  \end{bmatrix}
\end{equation}

Now we can write the forward map: $\mathbf{\Phi}(\mathbf{p}): \Omega \subset \mathbb{R}^5 \rightarrow \mathbb{R}^8$. More precisely,
\begin{equation}
  \Phi_i = \mathbf{H}_i(x,y,z) \cdot \hat{n}(\theta,\phi) 
\end{equation}
where $\mathbf{H}_i$ is the magnetic field calculated at $(x,y,z)$ due to coil $i$ using eqaution \eqref{eq:biot_savart_final}, and $\hat{n}(\theta,\phi)$ is the vector normal to the sensor.

Since this maps from a subset of $\mathbb{R}^5$ to $\mathbb{R}^8$ it is clear that the forward map is not surjective.

\subsection{The Jacobian}
Now that we have established that $\mathbf{\Phi}$ is not surjective we must move our attention to thinking about local invertibility rather than global. To do so we define the Jacobian:
\begin{equation}
  J(\mathbf{p}) = \frac{\partial \mathbf{\Phi}}{\partial \mathbf{p}} \in \mathbb{R}^{8\times5} \text{, } J_{i,j} = \frac{\partial \Phi_i}{\partial p_j}
  \label{eq:}
\end{equation}
Since $\mathbf{H}$ is a finite sum of Biot Savart applications it is a smooth map away from the current filaments. $\hat{n}$ is also a smooth map, therefore $\Phi$ will also be smooth away from the current filaments and the Jacobian will be well defined there.

A closed form for $\mathbf{J}$ is cumbersome find so we instead use finite differences to approximate $\mathbf{J}$

\begin{equation}
  \mathbf{J_{:,i}} \approx \frac{\mathbf{\Phi}\left(\mathbf{p} + h * \hat{e_i} \right) - \mathbf{\Phi(\mathbf{p})}}{h}
  \label{eq:}
\end{equation}
where $\hat{e}_i$ is the unit vector in the $i$th direction, and $h$ is a small real value ($h = 10^{-5}$ in our simulations).


\end{document}
